\section*{Abstract}

Predicting antibiotic resistance from a bacterial genome is far faster than testing for it in a laboratory, but a high score does not tell us how the model will behave on bacteria it has never seen. This thesis therefore examines two questions together: how genomic information should be represented, and how performance should be measured. Resistance arises in two ways: from mutations in genes the bacterium already carries, as \textit{gyrA} and \textit{parC} do for ciprofloxacin, or from genes acquired from elsewhere, as a TEM-1 beta-lactamase does for ampicillin. No single representation captures both, and existing multi-view methods combine views with fixed weights. We therefore built TriLens-MoE-Stack, which reads three views of each genome: 535{,}439 SNPs grouped by gene, a 32{,}441-gene presence profile, and four summary statistics per gene. A neural lens reads each view, a mixture-of-experts gate weights the three for every strain and antibiotic, and five classical learners are stacked on top. Because bacteria reproduce clonally, close relatives are common: 6{,}487 of our 6{,}962 \textit{Escherichia coli} fall into 458 sequence types, so a random split leaves relatives on both sides. We therefore held whole sequence types out of training. Of the fourteen predictive studies we reviewed, only two used population structure for their main split, and two more added a structure-aware check beside their headline result. Yet none corrects for the fact that a lineage holdout also changes the test prevalence, which directly affects PR-AUC. Our model achieved an average PR-AUC of 0.849 under a random split, compared with 0.803 for the gradient-boosting baseline, though its margin over the strongest single classical learner was only 0.009 and lies inside the bootstrap interval; under the sequence-type holdout it fell to 0.634. Matching the holdout prevalence to the training prevalence over 2{,}000 resampling rounds, with no retraining and no resistant isolate discarded, recovered the score to 0.723, so the prevalence-adjusted gap was 0.126 rather than the raw 0.215. The seven antibiotics then fell into two groups: ampicillin, ciprofloxacin and trimethoprim-sulfamethoxazole lost only 0.017 to 0.031, while cefotaxime, ceftazidime, amoxicillin-clavulanate and gentamicin lost 0.176 to 0.244, with nothing in between. With only seven drugs this is suggestive rather than conclusive, but the two groups differ in how their resistance is carried. \textit{gyrA} and \textit{parC} emerged as the top markers for ciprofloxacin, matching an independent published analysis of the same strains, whereas the larger losses point to a greater reliance on lineage. Attention weights, by contrast, were close to uniform, with entropies within ten per cent of the 8.52 ceiling for five of the seven antibiotics, and recovered only one of the eight known resistance genes that post-hoc attribution found. Retrained on 809 independent Giessen isolates, the classical component kept about 90\% of its Arcadia performance. No antibiotic met the applicable CLSI error limits, which we report as an unmet requirement. The contribution is therefore a protocol rather than a score: genomic resistance prediction should be evaluated with a lineage-aware split together with the prevalence correction it makes necessary, because generalisation depends on the resistance mechanism of the individual antibiotic and not on architecture alone.

\vspace{0.35cm}
\textbf{Keywords:} Antimicrobial resistance, genomic prediction, \textit{Escherichia coli}, lineage holdout, prevalence correction, hybrid learning, mixture-of-experts, multi-view learning
\pagebreak